\section{Machine-Learning Modelling of Scattering-Parameter Responses}
\label{sec:related-work}

% 1. From PCB datasets to direct S-parameter prediction
Schierholz \textit{et al.} introduced an open database of PCB interconnects and
demonstrated neural-network prediction of selected S-parameter magnitude and
phase \cite{schierholz2021sipi}. Their results established the feasibility of
data-driven S-parameter modelling and reproduced important response features,
but accuracy deteriorated for electrically long links with stronger
frequency-dependent behaviour. This highlights a central difficulty for
broadband surrogate modelling: more complex responses are harder to represent
accurately.

% 2 — Learning the broadband frequency response
Broadband prediction also raises an architectural challenge because neighbouring
frequency points are strongly correlated. Swaminathan \textit{et al.} identified
poor scaling of conventional fully connected networks for such responses
\cite{swaminathan2020demystifying}. Torun \textit{et al.} addressed this with
the spectral transposed convolutional neural network (S-TCNN), using transposed
convolution to generate a complete spectral response and exploit correlation
along frequency \cite{torun2019spectral}. This provides a more structured
alternative to treating frequency points independently, although the
architecture itself does not guarantee physically valid predictions.

% 3. Physical consistency, complex prediction, and the gap
Physical consistency is therefore an additional requirement. Chen and Tan
showed that an unconstrained neural network could produce non-physical
insertion-loss predictions, demonstrating that prediction error alone is
insufficient for evaluating an SI surrogate \cite{chen2021physics}. Akinwande
\textit{et al.} extended this principle to broadband complex $S_{11}$ and
$S_{21}$ using a complex-valued neural network with passivity and causality
enforcement \cite{akinwande2024complex}. This preserves phase information and
addresses physical behaviour explicitly, but still targets selected responses
rather than the complete multiport matrix. Taken together, the literature
progresses from selected S-parameter prediction towards structured, complex and
physically consistent broadband modelling; complete prediction of a
high-port-count complex S-matrix remains comparatively underexplored. This
motivates the reciprocal full 12-port formulation evaluated in this work.